\input ../preamble

\def\printmcount{\the\counter{section}.\the\counter{math counter}}
\setcounter{section}{1}

\centerline{\setfontandscale{bf}{20pt}Algebraic Geometry}
\medskip
\centerline{\bf Exercise \the\counter{section}}
\smallskip
\centerline{\it Ari Feiglin}

\bprob[title={Hartshorne II.1.8}]

    For an open subset $U\subseteq X$, consider the local sections functor $\Gamma(U,\bul)\colon\Ab_X\to\Ab$.
    Show that this functor is left exact.

\eprob

Being left exact is equivalent to being additive and preserving kernels.
We know that the presheaf sum of two sheaves and the presheaf kernel of a morphism of sheaves is a sheaf.
Thus $\Gamma(U,\bul)$ being left exact is due to the fact that it preserves presheaf sums and kernels, and when restricting our attention to sheaves
these limits are the same.
The local sections functor preserves limits (and colimits) in the presheaf category because limits in a functor category are taken componentwise.

\bprob[title={Hartshorne II.1.15}]

    Let $\F,\G$ be Abelian sheaves on $X$.
    Define a presheaf on $X$ by $U\mapsto\hom(\F|_U,\G|_U)$; show this defines a sheaf on $X$, called the {\emphcolor sheaf-hom}.

\eprop

We need to show that if $U=\bigcup_iU_i$ is an open cover and $\phi_i\colon\F|_{U_i}\to\G|_{U_i}$ are sheaf morphisms such that $\phi_i|_{U_i\cap U_j}=\phi_j|_{U_i\cap U_j}$,
then there is a $\phi\colon\F|_U\to\G|_U$ which restricts to $\phi_i$ on $U_i$.
Let $V\subseteq U$ be open, and let $s\in\F(V)$, then consider the sections $\phi_i(s|_{V\cap U_i})\in\G(V\cap U_i)$ of the open cover $V=\bigcup_i(V\cap U_i)$.
These sections agree on their common domains because $\phi_i$ do:
$$ \phi_i(s|_{V\cap U_i})|_{V\cap U_i\cap U_j} = \phi_i(s|_{V\cap U_i\cap U_j}) = \phi_j(s|_{V\cap U_i\cap U_j}) = \phi_j(s|_{V\cap U_j})|_{U_i} . $$
Because $\G$ is a sheaf there is a unique $\phi(s)\in\G(V)$ which restricts to $\phi(s|_{V\cap U_i})$ for each $i$.

This defines the $V$-component of a morphism of set-valued sheaves $\phi\colon\F|_U\to\G|_U$.
Indeed if $V\subseteq W\subseteq U$ are open, then for $s\in\F(W)$ we have that $\phi(s|_V)$ is the gluing of $\phi_i(s|_{V\cap U_i})$.
On the other hand $\phi(s)$ is the gluing of $\phi_i(s|_{W\cap U_i})$, and restricting this to $V\subseteq W$ gives the gluing of $\phi_i(s|_{V\cap U_i})$,
so that $\phi(s|_V)=\phi(s)|_V$ as desired.

Furthermore each component of $\phi$ is a morphism of Abelian groups.
For $V\subseteq U$, let $s,t\in\F(V)$ be sections.
Then
$$ \phi(s+t)|_{V\cap U_i} = \phi_i((s+t)|_{V\cap U_i}) = \phi_i(s|_{V\cap U_i}) + \phi_i(t|_{V\cap U_i}) = (\phi(s) + \phi(t))|_{V\cap U_i} , $$
and as $\G$ is a sheaf this means $\phi(s+t)=\phi(s)+\phi(t)$ as they agree on an open cover of $V$.
Thus $\phi\colon\F|_U\to\G|_U$ is a morphism in $\Ab_U$.

Finally it restricts to $\phi_i$ on $U_i$ because for $V\subseteq U_i$ we have that $\phi_i|_{V\cap U_j}=\phi_j|_{V\cap U_j}$ for all $j$.
So let $s\in\F(V)$, we have that
$$ \phi(s)|_{V\cap U_j} = \phi_j(s)|_{V\cap U_j} = \phi_i(s)|_{V\cap U_j} , $$
and thus $\phi(s)=\phi_i(s)$ as desired, meaning $\phi|_{U_i}=\phi_i$.
So we have shown that the sheaf-hom of two sheaves (in fact we needed only for $\G$ to be a sheaf) is a sheaf.

\bprob[title={Hartshorne II.1.16}, name={flasque}]

    A sheaf $\F$ on $X$ is {\emphcolor flasque} if its restriction maps are all surjective.
    \benum
        \item Show that constant sheaves on an irreducible space are flasque.
        \item Show that if $0\to\F'\to\F\to\F''\to0$ is exact with $\F'$ flasque then for any open $U\subseteq X$,
        $0\to\F'(U)\to\F(U)\to\F''(U)\to0$ is exact.
        \item Show that if $0\to\F'\to\F\to\F''\to0$ is exact with $\F',\F$ flasque, then $\F''$ is flasque.
        \item Show that if $f\colon X\to Y$ is continuous and $\F\in\Sh(X)$ is flasque, then $f_*\F\in\Sh(Y)$ is flasque.
        \item Let $\F\in\Sh(X)$ be a sheaf, then we define a new sheaf $\G$ on $X$ called the {\emphcolor sheaf of discontinuous sections} by
        $$ \G(U) = \prod_{p\in U}\F_p . $$
        Show that $\G$ is flasque and there is a natural injection $\F\inj\G$.
    \eenum

\eprob

\benum
    \item The restriction maps of $\ul A$ to nonempty subsets are all the identity $\id_A$ which is surjective.
    And the restrictions to the empty subset is the zero map which is surjective.

    \item First we claim that the presheaf quotient $(\F/\F')(U)=\F(U)/\F'(U)$ defines a sheaf.
    First it is separated, as if $U=\bigcup_iU_i$ is an open cover and $[s]\in\F(U)/\F'(U)$ satisfies $[s]|_{U_i}=0$ for all $i$, then
    $s|_{U_i}\in\F'(U_i)$ for all $i$.
    Because $\F'$ is a sheaf, by the sheaf property $s\in\F'(U)$ and so $[s]=0$.

    Now suppose we have $[s_i]\in\F(U_i)/\F'(U_i)$ such that $[s_i]|_{U_i\cap U_j}=[s_j]|_{U_i\cap U_j}$ for all $i,j$.
    This means that $(s_i-s_j)|_{U_i\cap U_j}\in\F'(U_i\cap U_j)$ for all $i,j$.
    Then we want to find an $s\in\F(U)$ such that $[s]|_{U_i}=[s_i]$ for all $i$, i.e.\ $s|_{U_i}-s_i\in\F'(U_i)$.

    Define the poset
    $$ \displaylines{
        \Sigma = \set{(V,s_V)}[V\subseteq U\hbox{ open},\ s_V\in\F(V),\ s_V|_{V\cap U_i}-s_i|_{U_i\cap V}\in\F'(U_i\cap V)\hbox{ for all $i$}],\cr
        (V,s_V) \leq (W,s_W) \iff V\subseteq W\hbox{ and } s_W|_V=s_V .
    } $$
    This poset is nonempty as it contains the zero section on the empty set.
    Furthermore it satisfies the conditions of Zorn's lemma: if $(V_\alpha,s_\alpha)_\alpha$ is a chain in $\Sigma$ then by definition
    the sections agree on their intersections (as it is a chain), and thus they glue to a section $s_V\in\F(V)$ on $V=\bigcup_\alpha V_\alpha$.
    And because $[s_V]|_{V_\alpha\cap U_i}=[s_i]|_{V_\alpha\cap U_i}$ for all $i,\alpha$, by separatedness we have that $[s_V]_{V\cap U_i}=[s_i]_{V\cap U_i}$
    so that $(V,s_V)\in\Sigma$ is an upper bound for the chain.

    So let $(V,s_V)\in\Sigma$ be maximal.
    If $U=V$ we are finished, otherwise suppose $V\subsetneq U$.
    Then there is an $i$ such that $V\cap U_i\neq U_i$.
    By definition $(s_V-s_i)|_{V\cap U_i}=t\in\F'(U_i\cap V)$, and by flasqueness we may view $t$ as a section of $\F'(U_i)$ so $(s_V-s_i)|_{V\cap U_i}=t|_{V\cap U_i}$.
    Consider then the sections
    $$ s_V \in \F(V),\qquad s_i + t \in \F'(U_i) \subseteq \F(U_i) . $$
    These agree on the intersection $V\cap U_i$, and thus extend to a section $\bar s_V$ on $\bar V=V\cup U_i\supsetneq V$.
    Notice that for a $j$,
    $$ \displaylines{
        [\bar s_V]|_{V\cap U_j} = [s_V]|_{V\cap U_j} = [s_j]|_{V\cap U_j} ,\cr
        [\bar s_V]|_{U_i\cap U_j} = [s_i+t]|_{U_i\cap U_j} = [s_i]|_{U_i\cap U_j} = [s_j]|_{V\cap U_j}.
    } $$
    Thus $[\bar s_V]$ agrees with $[s_j]$ on an open cover of $\bar V$ and thus on all of $\bar V$.
    So we have shown that $(\bar V,\bar s_V)\in\Sigma$, but it is a proper extension of $(V,s_V)$ contradicting its maximality.
    Thus $U=V$ and we are finished.

    Given an exact sequence $0\to\F'\to\F\to\F''\to0$, then we have $\F''=\coker(\F'\to\F)$ which is just $\F/\F'$.
    As we showed, this is the presheaf quotient, whose local quotients are $\F(U)/\F'(U)$.
    So we get the desired result, but to be explicit we have an isomorphism of short exact sequences

    \diagram{
        0&\F'&\F&\F''&0\cr
        0&\F'&\F&\F/\F'&0
    }{
        \diagarrow{from={1,1}, to={1,2}}
        \diagarrow{from={1,2}, to={1,3}}
        \diagarrow{from={1,3}, to={1,4}}
        \diagarrow{from={1,4}, to={1,5}}
        \diagarrow{from={2,1}, to={2,2}}
        \diagarrow{from={2,2}, to={2,3}, text=\subseteq, y distance=.2cm}
        \diagarrow{from={2,3}, to={2,4}}
        \diagarrow{from={2,4}, to={2,5}, text=\pi, y distance=.2cm}
        \diagarrow{from={1,2}, to={2,2}, text=\cong, x distance=.2cm}
        \diagarrow{from={1,3}, to={2,3}, text=\cong, x distance=.2cm}
        \diagarrow{from={1,4}, to={2,4}, text=\cong, x distance=.2cm}
    }

    Applying $\Gamma(U,\bul)$ maps the bottom sequence to a short exact sequence, and the top sequence is isomorphic to it and thus is also exact.

    \item Let $U\subseteq V$ be open in $X$, then by (2) both of the rows in the following sequence are exact

    \diagram{
        0&\F'(V)&\F(V)&\F''(V)&0\cr
        0&\F(U)'&\F(U)&\F''(U)&0
    }{
        \diagarrow{from={1,1}, to={1,2}}
        \diagarrow{from={1,2}, to={1,3}}
        \diagarrow{from={1,3}, to={1,4}, right cap=>>}
        \diagarrow{from={1,4}, to={1,5}}
        \diagarrow{from={2,1}, to={2,2}}
        \diagarrow{from={2,2}, to={2,3}}
        \diagarrow{from={2,3}, to={2,4}, right cap=>>}
        \diagarrow{from={2,4}, to={2,5}}
        \diagarrow{from={1,2}, to={2,2}, right cap=>>}
        \diagarrow{from={1,3}, to={2,3}, right cap=>>}
        \diagarrow{from={1,4}, to={2,4}}
    }

    where the vertical arrows are the restriction maps.
    The last horizontal maps are surjective because the rows are exact, and the first two vertical maps are surjective because $\F',\F$ are flasque.
    In particular the composition $\F(V)\to\F''(V)\to\F''(U)$ is surjective, which implies that restriction $\F''(V)\to\F''(U)$ is.
    Thus $\F''$ is flasque, as desired.

    \item Let $U\subseteq V\subseteq Y$ be open.
    Then by definition
    $$ f_*\F(V) \xtos\res f_*\F(U) \quad=\quad \F(f^{-1}V) \xtos\res \F(f^{-1}U) , $$
    which is surjective as $\F$ is flasque.
    So $f_*\F$ is flasque as well.

    \item The restriction map between $U\subseteq V\subseteq X$ for $\G$ is
    $$ (s_p\in\F_p)_{p\in V} \in \G(V) \mapsto (s_p\in\F_p)_{p\in U} \in \G(U) $$
    is surjective as for $(s_p\in\F_p)_{p\in U}\in\G(U)$, we can define $(\bar s_p)_{p\in V}\in\G(V)$ by
    $$ \bar s_p = \cases{s_p&$p\in U$\cr 0&else} , $$
    and then clearly $(\bar s_p)_V$ restricts to $(s_p)_U$.

    Define
    $$ \F(U) \to \G(U),\qquad s\mapsto(s_p)_{p\in U} . $$
    This is an Abelian group morphism because its coordinates are the canonical projections $\F(U)\to\F_p$ which are Abelian morphisms.
    This is injective as if $s\in\F(U)$ is in the kernel, then $s_p=0$ for all $p\in U$.
    Thus for all $p\in U$ there is an open neighborhood $p\in U_p\subseteq U$ such that $s|_{U_p}=0$.
    The $U_p$ then cover $U$ and $s|_{U_p}=0$ for all $p$, and as $\F$ is separated this means $s=0$.
    So the kernel is trivial and thus the map is an injection.

    Furthermore these maps clearly assemble into a morphism of sheaves.
    Indeed, if $U\subseteq V$ then both paths $\F(V)\to\G(U)$ in the square
    \commsq{\F(V)&\G(V)\cr\F(U)&\G(U)}
        {}{}{}{}
    map $s\in\F(V)$ to $(s_p)_{p\in U}$ (recall that $(s|_U)_p=s_p$ for $p\in U$ and $s\in\F(V)$).
    The components of this morphism are injections and thus this map $\F\to\G$ is injective.
\eenum

\bprob[title={Hartshorne II.1.17}]

    Let $X$ be a topological space, $p\in X$, and $A\in\Ab$ be an Abelian group.
    Define the {\emphcolor skyscraper sheaf} by
    $$ p_*A(U) = \cases{A&$p\in U$\cr0&else} . $$
    Verify that this is a sheaf, and that its stalks satisfy
    $$ (p_*A)_q = \cases{A&$q\in\cl p$\cr0&else} . $$
    Also show that $p_*A$ can be viewed as the pushforward of the constant sheaf $A$ on $*$ via $*\to A$ which maps $*$ to $p$, or as
    the pushforward of the constant sheaf $A$ on the irreducible subspace $\cl p\subseteq A$ along its inclusion to $A$.

\eprob

We first show the last claim, which will automatically show that $p_*A$ is a sheaf.
Let $i_p\colon *\to A$ be the map $i_p(*)=p$.
Then $A$ is a constant sheaf on $*$ and its pushforward satisfies for $U\subseteq X$,
$$ i_{p,*}A(U) = A(i_p^{-1}U) = \cases{A&$i_p^{-1}(U)\neq\emptyset$\cr0&else} , $$
and because $i_p^{-1}U$ is empty iff $p\notin U$ we see that
$$ i_{p,*}A(U) = p_*A(U) . $$
The restriction maps of $i_{p,*}A$ are those of the constant sheaf $A$, i.e.\ the identities when $p\in U\subseteq V$ and zero otherwise,
which are the restriction maps of $p_*A$, so $p_*A=i_{p,*}A$ as desired.

Note that $q\in\cl p$ iff every neighborhood of $q$ contains $p$.
So if we denote by $i_p\colon\cl p\inj A$ the inclusion, then
$$ i_{p,*}A(U) = \cases{A&$i_p^{-1}U\neq\emptyset$\cr0&else} , $$
and $i_p^{-1}U$ is nonempty iff $\cl p\cap U$ is nonempty, which means there is a $q\in\cl p$ for which $U$ is a neighborhood.
This is iff $p\in U$, and so we see that $i_{p,*}A(U)=p_*A(U)$.
And again the restriction maps are equal.

Now for $q\in X$ we have that
$$ (p_*A)_q = \colim_{q\in U}p_*A(U) . $$
If $q\notin\cl p$ then there is a neighborhood $U$ of $q$ which does not contain $p$.
The set of all such neighborhoods is cofinal in $\open_q(X)^\op$ and thus we can restrict the colimit to this subposet:
$$ (p_*A)_q = \colim_{q\in U,p\notin U}p_*A(U) = \colim_U0 = 0 . $$
And if $q\in\cl p$ then all of $q$'s neighborhoods contain $p$, so the filtered colimit above is just the filtered colimit of the constant functor valued at $A$,
which is equal to $A$, as desired.

\bprob[title={Hartshorne II.1.19}]

    Let $Z\subseteq X$ be closed and $i\colon Z\inj X$ the inclusion.
    Define $U=X-Z$ to be its complement and $j\colon U\inj X$ its inclusion.
    \benum
        \item Let $\F$ be a sheaf on $Z$.
        Show that for $p\in X$ the stalk $(i_*\F)_p$ is $\F_p$ if $p\in Z$ and $0$ otherwise.
        \item Let $\F$ be a sheaf on $U$, and define $j_!\F$ to be the sheaf on $X$ associated to the presheaf
        $$ j_!^\pre\F(V) = \cases{\F(V)&if $V\subseteq U$\cr 0&else} . $$
        Show that for $p\in X$ the stalk of this sheaf satisfies
        $$ (j_!\F)_p = \cases{\F_p&$p\in U$\cr0&else}, $$
        and that $j_!\F$ is uniquely determined by this property and that $(j_!\F)|_U=\F$.
        The sheaf $j_!\F$ is called the {\emphcolor extension of $\F$ by zero}.
        \item Now let $\F$ be a sheaf; show that there is an exact sequence of sheaves on $X$
        $$ 0 \to j_!(\F|_U) \to \F \to i_*(\F|_Z) \to 0 . $$
    \eenum

\eprob

\benum
    \item Note that the sheaf $i_*\F$ is defined by, for $V\subseteq X$ open
    $$ (i_*\F)(V) = \F(i^{-1}V) = \F(V\cap Z) . $$
    Thus for $p\in X$,
    $$ (i_*\F)_p = \colim_{p\in V}\F(V\cap Z) . $$
    If $p\in Z$ then all of $p$'s neighborhoods in $Z$ are of the form $V\cap Z$ for $V\subseteq X$ open and so this is equal to
    $$ (i_*\F)_p = \colim_{p\in W\subseteq Z}\F(W) = \F_p $$
    (where the last equality is by definition).
    Otherwise if $p\notin Z$ then as $Z$ is closed there is a neighborhood $p\in W$ disjoint from $Z$.
    Then $W\cap Z=\emptyset$ is terminal in the category over which we take the colimit (which we view as $\set{V\cap Z}_V$ ordered by reverse inclusion).
    The colimit over a category with a terminal object is equal to the functor at that object so
    $$ (i_*\F)_p = \colim_{p\in V}\F(V\cap Z) = \F(W\cap Z) = \F(\emptyset) = 0 . $$

    \item We know that the stalks of the sheaf associated to a presheaf are isomorphic to the stalks of the presheaf.
    This is because we have adjoints
    $$ \Psh(X) \adjarrow{\sh}{\inc} \Sh(X) \adjarrow{\bul_p}{p_*\bul} \Ab , $$
    and so we have adjoints
    $$ \Psh(X) \adjarrow{\bul^\sh_p}{p_*\bul} \Ab . $$
    But taking stalks is left-adjoint to the skyscraper sheaf over the category of presheaves.
    Thus $\bul^\sh_p$ and $\bul_p$ are both left-adjoints of $\Ab\xtos{p_*\bul}\Psh(X)$, and thus isomorphic.
    So we can focus on the stalks of $j_!^\pre\F$; these are isomorphic to the stalks of $j_!\F$.

    As before for $p\in X$,
    $$ (j_!^\pre\F)_p = \colim_{p\in V}j_!^\pre\F(V) . $$
    If $p\in U$ then the subposet of neighborhoods of $p$ contained in $U$ is cofinal in $\open_p(X)^\op$ (the opposite poset of $p$'s neighborhoods) and thus
    we can restrict the colimit to this category:
    $$ (j_!^\pre\F)_p = \colim_{p\in V\subseteq U}j_!^\pre\F(V) = \colim_{p\in V\subseteq U}\F(V) = \F_p . $$
    Conversely if $p\notin U$ then no neighborhood of $p$ is contained in $U$ (as it is open) and thus
    $$ (j_!^\pre\F)_p = \colim_{p\in V}j_!^\pre\F(V) = \colim_{p\in V}0 = 0 $$
    as the filtered colimit of a constant functor is the constant.

    If $\G$ is another sheaf satisfying $\G|_U=\F$ and $\G_p=\F_p$ if $p\in U$ and $0$ otherwise, then $\G$ is isomorphic to $j_!\F$.
    Indeed, define the morphism $\phi\colon\G\to j_!^\pre\F$ whose component for $V\subseteq X$ is $\id$ if $V\subseteq0$ and $0$ otherwise.
    This is clearly a morphism of presheaves, and thus composition with the sheafification map $\G\to j_!^\pre\F\to j_!\F$, which we will also denote by $\phi$,
    is a morphism of sheaves.
    For every $p\in X$, $\phi_p$ is an isomorphism.
    On $U$ $\phi$ is the identity and so this is clear, and for $p\notin U$ then $\G_p,(j_!\F)_p$ are both zero and thus $\phi_p=0$ is automatically an isomorphism.
    Because $\phi_p$ is an isomorphism for all $p\in X$, $\phi$ is an isomorphism $\G\cong j_!\F$.

    \item Define $j_!^\pre(\F|_U)\to\F$ to be the obvious map, whose component on $V$ is either $\id$ or $0$ depending on if $V\subseteq U$ or not.
    By the universal property of sheafification this defines a map $j_!(\F|_U)\to\F$.
    The stalks of this map are the same as the stalks of $j_!^\pre(\F|_U)\to\F$ (by the argument in the previous point), i.e.\ the identity when $p\in U$
    and zero otherwise.

    Define $\F\to i_*(\F|_Z)$ to be the unit of the adjunction $\bul|_Z=i^{-1}\dashv i_*$.
    Then for $p\in Z$ we have by (1)
    $$ i_*(\F|_Z)_p = (\F|_Z)_p = \F_p . $$
    The last equality is because $i^{-1}\dashv i_*$ and $\bul_p\dashv p_*$ are adjunctions and thus $(i^{-1}\bul)_p\dashv i_*(p_*\bul)$.
    Clearly $f_*g_*\bul=(f\circ g)_*\bul$ and thus $i_*(p_*\bul)=i(p)_*\bul$, which has left adjoint $\bul_{i(p)}$, so that $(i^{-1}\bul)_p=\bul_{i(p)}$ (up to natural isomorphism).
    And for $p\in U=X-Z$ we have by (1)
    $$ i_*(\F|_Z)_p = 0 . $$

    Consider the sequence of these maps
    $$ 0 \to j_!(\F|_U) \to \F \to i_*(\F|_Z) \to 0 . $$
    Then taking the stalk at $p\in Z$ we get the sequence
    $$ 0 \to 0 \to \F_p \to i_*(\F|_Z)_p = \F_p \to 0 $$
    which is exact as the non-trivial arrow is an isomorphism.
    And for $p\notin Z$ we get the sequence
    $$ 0 \to j_!(\F|_U)_p = \F_p \to \F_p \to 0 \to 0 , $$
    which is again exact.
    So taking stalks at all $p\in X$ gives an exact sequence, so the original sequence is exact.
\eenum

\bprob[title={Hartshorne II.1.20}]

    Let $Z\subseteq X$ be a closed subset and $\F$ a sheaf on $X$.
    Define $\Gamma_Z(X,\F)$ to be the subgroup of $\Gamma(X,\F)$ consisting of sections whose supports are contained in $Z$.
    \benum
        \item Show that the presheaf $V\mapsto\Gamma_{Z\cap V}(V,\F|_V)$ is a sheaf.
        It is called the {\emphcolor subsheaf of $\F$ with supports in $Z$} and is denoted $\H^0_Z(\F)$.
        \item Let $U=X-Z$ and $j\colon U\to X$ the inclusion.
        Show that there is an exact sequence of sheaves
        $$ 0 \to \H^0_Z(\F) \to \F \to j_*(\F|_U) $$
        and if $\F$ is flasque the last arrow is surjective.
    \eenum

\eprob

\benum
    \item Let $U=\bigcup_iU_i$ be an open cover and $s_i\in\Gamma_{Z\cap U_i}(U_i,\F_{U_i})$ such that $s_i|_{U_i\cap U_j}=s_j|_{U_i\cap U_j}$.
    Then by gluing, these form a unique section $s\in\Gamma(U,\F|_U)$ which restricts to $s_i$ on $U_i$.
    Note that for $p\in U_i$, $s_p=s_{i,p}$ and thus $p\in\supp s$ iff $p\in\supp s_i$ so that
    $$ \supp s = \bigcup_i\supp s_i \subseteq \bigcup_i(Z\cap U_i) = Z\cap\bigcup_iU_i = Z\cap U , $$
    meaning $s\in\Gamma_{Z\cap U}(U,\F|_U)$ as desired.
    Uniqueness of gluing follows from the uniqueness of gluing in $\F$.

    \item The inclusion $\H^0_Z(\F)\subseteq\F$ gives the first arrow in the sequence.
    The map $\F\to j_*(\F|_U)$ is given by restriction $\F(V)\to \F(U\cap V)$.
    Because $\H^0_Z(\F)$ is a sheaf, the image of $\H^0_Z(\F)\inj\F$ {\it is\/} $\H^0_Z(\F)$ (the image of a sheaf morphism is the sheafification of the presheaf image).
    So we just need to show that $\ker(\F\to j_*(\F|_U))=\H^0_Z(\F)$.

    This is equivalent to the claim that for $s\in\F(V)$, $s|_{U\cap V}=0$ iff $s\in\Gamma_{Z\cap V}(V,\F)$.
    Note that $s\in\Gamma_{Z\cap V}(V,\F)$ iff $s_p=0$ when $p\in V-Z=U\cap V$.
    And $s|_{U\cap V}=0$ iff $s_p=0$ for all $p\in U\cap V$, proving the desired claim.
    This last equivalence is due to \refmath{flasque} (5), as $\F(U)\to\prod_{p\in U}\F_p$ is injective.

    Now if $\F$ is flasque then $\F\to j_*(\F|_U)$ is clearly surjective, because its components are the surjections $\F(V)\to\F(U\cap V)$.
    So this map is surjective on each component, thus surjective as a presheaf morphism and in particular as a sheaf morphism (an epimorphism in a category is an epimorphism in any subcategory
    containing it).
\eenum

\bprob[title={Hartshorne II.1.22}]

    Let $X$ be a topological space with open cover $X=\bigcup_iU_i$.
    For each $i$ let $\F_i$ be a sheaf on $U_i$, and suppose for each $i,j$ we have an isomorphism $\phi_{ij}\colon\F_i|_{U_i\cap U_j}\to\F_j|_{U_i\cap U_j}$.
    Finally suppose that these isomorphisms satisfy the {\emphcolor cocycle conditions}:
    \benum
        \item $\phi_{ii}=\id$
        \item for $i,j,k$, $\phi_{ik}=\phi_{jk}\circ\phi_{ij}$ on $U_i\cap U_j\cap U_k$.
    \eenum
    Show that there exists a unique sheaf $\F$ on $X$ with isomorphisms $\psi_i\colon\F|_{U_i}\to\F_i$ such that $\psi_j=\phi_{ij}\circ\psi_i$ on $U_i\cap U_j$.

\eprob

Let us write $U_{i_0,\dots,i_n}$ for the intersection $U_{i_0}\cap\cdots\cap U_{i_n}$.
For $V\subseteq X$ define
$$ \F(V) = \set{(s_i)_i\in\prod_i\F_i(V\cap U_i)}[s_i|_{V\cap U_{ij}}=s_j|_{V\cap U_{ij}}\hbox{ via $\phi_{ij}$, for all $i,j$}] . $$
By $x=y$ via $\phi$ we mean $\phi(x)=y$.
This naturally carries the structure of an Abelian group by $(s_i)_i+(t_i)_i=(s_i+t_i)_i$.
Then for $V\subseteq W$ we define $\F(W)\to\F(V)$ by
$$ (s_i)_i \in \F(W) \mapsto (s_i|_{V\cap U_i}) \in \F(V) . $$
This is clearly well-defined and natural (i.e.\ makes $\F$ into a presheaf).
Furthermore $\F$ is a sheaf.
Let $V=\bigcup_\alpha V_\alpha$ be an open cover and $(s_i^\alpha)\in\F(V_\alpha)$ be sections such that $(s_i^\alpha)|_{V_\alpha\cap V_\beta}=(s_i^\beta)_{V_\alpha\cap V_\beta}$
for all $\alpha,\beta$.
In particular this means $s_i^\alpha$ and $s_i^\beta$ agree with each other on $V_{\alpha,\beta}\cap U_i$.
Because $\F_i$ is a sheaf, this means they uniquely glue to give $s_i\in\F(V\cap U_i)$ (by uniqueness of $\F_i$'s gluing this is unique).
Now $(s_i)_i\in\F(V)$ because for each $i,j$,
$$ \phi_{ij}(s_i|_{U_{ij}})|_{V_\alpha} = \phi_{ij}(s_i|_{V_\alpha\cap U_{ij}}) = s_j|_{V_\alpha\cap U_{ij}}\qquad\hbox{for all $\alpha$} . $$
And thus because $\F_i$ is separated this means $\phi_{ij}(s_i|_{U_{ij}})=s_j|_{U_{ij}}$ as desired.

Now for $V\subseteq U_i$ define
$$ \psi_i\colon\F(V) \to \F_i(V),\qquad (s_j\in\F_i(V\cap U_j))_j \mapsto s_i \in \F_i(V) . $$
This is clearly a well-defined morphism.
Further define
$$ \tilde\psi_i\colon\F_i(V)\to\F(V),\qquad s \mapsto (\phi_{ij}(s|_{V\cap U_j}))_j . $$
This is well-defined: for $j,k$ we have
$$ \phi_{jk}(\phi_{ij}(s|_{V\cap U_j})|_{V\cap U_{jk}}) = \phi_{jk}\circ\phi_{ij}(s|_{V\cap U_{ijk}}) = \phi_{ik}(s|_{V\cap U_{ijk}}) = \phi_{ik}(s|_{V\cap U_k})|_{V\cap U_{jk}} . $$
Furthermore these clearly assemble into a morphism $\F_i\to\F|_{U_i}$.
Now $\psi_i,\tilde\psi_i$ are inverses: for $s\in\F_i(V)$,
$$ \psi_i\circ\tilde\psi_i(s) = \psi_i(\phi_{ij}(s|_{V\cap U_j}))_j = \phi_{ii}(s|_{V\cap U_i}) = s $$
by the cocycle condition (1).
And for $(s_j)_j\in\F(V)$,
$$ \tilde\psi_i\circ\psi_i(s_j)_j = \tilde\psi_i(s_i) = (\phi_{ij}(s_i|_{V\cap U_j}))_j . $$
So we want to show that $\phi_{ij}(s_i|_{V\cap U_j})=s_j$.
Because $(s_j)_j\in\F(V)$ we have $\phi_{ij}(s_i|_{V\cap U_j})=s_j|_{V\cap U_j}=s_j$ as desired.

And for $V\subseteq U_i\cap U_j$ and $(s_k)_k\in\F(V)$ we have that
$$ \phi_{ij}\circ\psi_i(s_k)_k = \phi_{ij}(s_i),\qquad \psi_j(s_k)_k = s_j . $$
And as before because $(s_k)_k\in\F(V)$ we have $\phi_{ij}(s_i)=s_j$ on $V\cap U_i\cap U_j=V$, so $\phi_{ij}\circ\psi_i=\psi_j$.

The sheaf $\F$ satisfies the universal property:
$$ \hom(\F,\G) \cong \set{(\rho_i\colon\F_i\to\G|_{U_i})_i}[\rho_i=\rho_j\phi_{ij}\hbox{ on $U_{ij}$}] . \eqcnt $$
Clearly every morphism $\rho\colon\F\to\G$ defines
$$ \rho_i\colon \F_i \xtos{\psi_i^{-1}} \F|_{U_i} \xtos{\rho|_{U_i}} \G|_{U_i} $$
and on $U_{ij}$,
$$ \rho_j\circ\phi_{ij} = \rho|_{U_{ij}}\circ\psi_j^{-1}\circ\phi_{ij} = \rho|_{U_{ij}}\circ\psi_i^{-1} = \rho_i , $$
as desired.
And given $(\rho_i)_i$ define $\rho\colon\F(V)\to\G(V)$ by mapping $(s_i\in\F_j(V\cap U_i))_i$ to the gluing of $\rho_i(s_i)\in\G(V\cap U_i)$.
By definition these sections agree on their intersections:
$$ \rho_j(s_j)|_{U_{ij}} = \rho_j(\phi_{ij}(s_i|_{U_{ij}})) = \rho_i(s_i)|_{U_{ij}} , $$
so such a gluing exists and is unique.

These two processes are mutual inverses, and thus \gotoleqn{} is a natural isomorphism.
Because universal properties define objects up to unique isomorphism, we are finished.

\bye

